% DERIVED from formal_layer.json — read-only, do not edit.
\begin{definitionv}{def:holder-norm}[Coordinate-partial Hölder norm]
The extended-real Hölder norm \(\|f\|_{H^s}\), for \(f:\mathbb R^d\to\mathbb R\), \(d\in\{0,1,\ldots\}\), and \(s\in\mathbb R\), is
\[
\|f\|_{H^s}
=
\begin{cases}
\displaystyle
\max\left\{
\sup_{\substack{\kappa\in\{0,1,\ldots\}^d,\;
|\kappa|\le \max\{\lceil s\rceil-1,0\}\\
x\in[0,1]^d}}
|\partial^\kappa f(x)|,\;
\sup_{\substack{\kappa\in\{0,1,\ldots\}^d,\;
|\kappa|=\max\{\lceil s\rceil-1,0\}\\
x,x'\in[0,1]^d,\;x\ne x'}}
\frac{|\partial^\kappa f(x)-\partial^\kappa f(x')|}
{\|x-x'\|_2^{\,s-\max\{\lceil s\rceil-1,0\}}}
\right\},
& f\in C^{\max\{\lceil s\rceil-1,0\}}([0,1]^d),\\[1ex]
+\infty,
& \text{otherwise}.
\end{cases}
\]
Here \(|\kappa|=\sum_{i=1}^d\kappa_i\), and \(\partial^\kappa f(x)\) is the iterated derivative within \([0,1]^d\), evaluated on coordinate directions with coordinate \(i\) repeated \(\kappa_i\) times. The differentiability condition means continuous differentiability of the indicated order within the cube. Both suprema are taken in the extended nonnegative reals, with an empty supremum equal to zero.
\end{definitionv}

\begin{definitionv}{synth_5}[Observed response selection]
The observed response \(Y\) for a full record \((X,A,Y(0),Y(1))\), where \(X\in\mathbb R^d\) is the covariate vector for a nonnegative integer \(d\), \(A\in\{0,1\}\) is the treatment indicator, and \(Y(0),Y(1)\in\{0,1\}\) are the potential responses under control and treatment, is defined by
\[
Y=
\begin{cases}
Y(0), & A=0,\\
Y(1), & A=1.
\end{cases}
\]
The corresponding observed record is \((X,A,Y)\).
\end{definitionv}

\begin{assumptionv}{ass:uniform-design}[Uniform covariate design]
The covariate vector \(X\) satisfies
\[
\mathcal L_P(X)=\mathrm{Unif}([0,1]^d).
\]
\end{assumptionv}

\begin{assumptionv}{ass:exchangeability}[Conditional treatment exchangeability]
The treatment indicator \(A\) and the potential responses \(Y(0),Y(1)\) satisfy
\[
A\perp(Y(0),Y(1))\mid X,
\]
where \(X\) is the covariate vector.
\end{assumptionv}

\begin{assumptionv}{ass:overlap}[Pointwise overlap on the cube]
The overlap condition for a primitive law \(P\), at the overlap level \(\varepsilon\in\mathbb R\), is \(0<\varepsilon\) together with the requirement that the raw Borel propensity witness carried by \(P\) takes values in \([\varepsilon,1-\varepsilon]\) at every \(x\in[0,1]^d\).
\end{assumptionv}

\begin{assumptionv}{ass:propensity-holder}[Hölder smoothness of the propensity]
The designated Borel propensity witness \(e_P\) satisfies
\[
\|e_P\|_{H^\alpha}\le L.
\]
\end{assumptionv}

\begin{assumptionv}{ass:control-holder}[Hölder smoothness of the control mean]
The designated Borel control-mean witness \(\mu_{0,P}\) satisfies
\[
\|\mu_{0,P}\|_{H^\beta}\le L.
\]
\end{assumptionv}

\begin{assumptionv}{ass:effect-holder}[Hölder smoothness of the causal contrast]
The canonical continuous causal contrast \(\tau_P\) satisfies
\[
\|\tau_P\|_{H^\gamma}\le L.
\]
\end{assumptionv}

\begin{assumptionv}{ass:control-interior}[Control response bounds]
The designated Borel control-mean witness \(\mu_{0,P}\) satisfies
\[
\mu_{0,P}(x)\in[0,1]
\qquad\text{for every }x\in[0,1]^d.
\]
\end{assumptionv}

\begin{assumptionv}{ass:treated-interior}[Treated response bounds]
The designated Borel treated-mean witness \(\mu_{1,P}\) satisfies
\[
\mu_{1,P}(x)\in[0,1]
\qquad\text{for every }x\in[0,1]^d.
\]
\end{assumptionv}

\begin{definitionv}{def:primitive-class}[Primitive law class \(\mathcal M(d,\alpha,\beta,\gamma,L,\varepsilon)\)]
\(\mathcal M(d,\alpha,\beta,\gamma,L,\varepsilon)\) is the class of primitive Borel laws \(P=\mathcal L(X,A,Y(0),Y(1))\), where \(X\) is the covariate vector, \(A\) is the binary treatment indicator, and \(Y(0)\) and \(Y(1)\) are the binary potential responses under control and treatment. For public parameters \(L>1/2\) and \(\varepsilon\in(0,1/2)\), membership is defined by the following properties:
\begin{itemize}
\item \textbf{(Uniform design.)} The law satisfies \cref{obj:ass:uniform-design}.
\item \textbf{(Exchangeability.)} The law satisfies \cref{obj:ass:exchangeability}.
\item \textbf{(Conditional means.)} There exist Borel functions
\[
g_e,g_0,g_1:[0,1]^d\to[0,1]
\]
such that
\[
\begin{aligned}
g_e(X)&=\mathbb E_P[A\mid X] &&P\text{-almost surely},\\
g_0(X)&=\mathbb E_P[Y(0)\mid X] &&P\text{-almost surely},\\
g_1(X)&=\mathbb E_P[Y(1)\mid X] &&P\text{-almost surely}.
\end{aligned}
\]
\item \textbf{(Overlap and smoothness.)} These functions satisfy \cref{obj:ass:overlap,obj:ass:propensity-holder,obj:ass:control-holder,obj:ass:effect-holder}, with \(g_e\), \(g_0\), and \(g_1-g_0\) as the propensity, control mean, and causal contrast.
\item \textbf{(Response bounds.)} The arm means satisfy \cref{obj:ass:control-interior,obj:ass:treated-interior}.
\end{itemize}
For each admitted law, \(e_P,\mu_{0,P},\mu_{1,P}\) designate admissible functions \(g_e,g_0,g_1\), respectively. The function \(\tau_P\) is the canonical continuous representative of the almost-sure conditional causal contrast, and the designated functions satisfy
\[
\mu_{1,P}(x)-\mu_{0,P}(x)=\tau_P(x)
\qquad\text{for every }x\in[0,1]^d.
\]
\end{definitionv}

\begin{assumptionv}{ass:sample-law}[Original experiment sampling law]
The original dataset \(\mathcal D_{n,m}\) and independent uniform randomizer \(U\) have joint law
\[
\mathcal L_P(\mathcal D_{n,m},U)=\mathsf E_{n,m}(P),
\]
where \(\mathsf E_{n,m}(P)\) is the product sampling law of the original dataset and randomizer.
\end{assumptionv}

\begin{definitionv}{def:risk}[Absolute-error risk $\mathcal R_{n,m}(T,P)$]
The absolute-error risk is
\[
\mathcal R_{n,m}(T,P)
=
\int
\left|T(\mathcal D_{n,m},U)-\tau_P(x_0)\right|
\,d\mathsf E_{n,m}(P).
\]
Here \(T\) is a Borel real-valued decision using the original dataset \(\mathcal D_{n,m}\) and independent uniform randomizer \(U\), \(x_0\) is the fixed interior covariate profile, and \(\mathsf E_{n,m}(P)\) is the product sampling law of the original dataset and randomizer. The target \(\tau_P(x_0)\) is the canonical conditional causal contrast from \cref{obj:def:primitive-class}.
\end{definitionv}

\begin{definitionv}{def:minimax}[Minimax risk \(R(n,m)\)]
\[
R(n,m)=\inf_{T\ {\rm Borel}}\sup_{P\in\mathcal M}\mathcal R_{n,m}(T,P).
\]
Here \(\mathcal M\) is the primitive law class at the fixed public parameters, and \(\mathcal R_{n,m}(T,P)\) is the expected absolute error for the point contrast under law \(P\). The infimum ranges over Borel real-valued decisions \(T\) using the original dataset of \(n\) labeled records and \(m\) independent outcome-free treatment records, together with an independent uniform randomizer.
\end{definitionv}

\begin{definitionv}{def:oracle-risk}[Supplied propensity minimax risk]
The supplied propensity minimax risk \(R^e(n,m)\), for \(d,n,m\in\mathbb N\) and \(\alpha,\beta,\gamma,L,\varepsilon\in\mathbb R\), is the extended nonnegative quantity
\[
R^e(n,m)=\inf_{T^e}\sup_{P\in\mathcal M}
\int
\left|T^e(\mathcal D_{n,m},U,e_P)-\tau_P(x_0)\right|
\,d\mathsf E_{n,m}(P).
\]
Here \(\mathcal M\) is the primitive law class at these parameters, and \(e_P\) and \(\tau_P\) are its designated admissible propensity function and canonical causal contrast, as defined in \cref{obj:def:primitive-class}. The evaluation profile is \(x_0=(1/2,\ldots,1/2)\). The original dataset \(\mathcal D_{n,m}\) comprises \(n\) labeled records and \(m\) independent outcome-free treatment records; \(U\) is the independent uniform randomizer on \([0,1]\), and \(\mathsf E_{n,m}(P)\) is their product sampling law from \cref{obj:ass:sample-law}. The infimum ranges over real-valued decisions \(T^e\) receiving the dataset, randomizer, and a supplied function from the covariate space to \(\mathbb R\), with Borel measurability in \((\mathcal D_{n,m},U)\) for every fixed supplied function.
\end{definitionv}

\begin{definitionv}{synth_3}[Localization cube and uniform localization measure]
For an integer \(d\ge1\), let \(x_0=(1/2,\ldots,1/2)\in[0,1]^d\). For \(0<h\le1/2\), define the localization cube of side length \(h\), centered at \(x_0\), by
\[
C_h=x_0+[-h/2,h/2]^d.
\]
The uniform probability measure \(\nu_h\) on \(C_h\) has density \(w_h\) with respect to \(d\)-dimensional Lebesgue measure \(dx\) on \([0,1]^d\):
\[
w_h(x)=h^{-d}\mathbf1\{x\in C_h\},
\qquad
d\nu_h(x)=w_h(x)\,dx.
\]
Equivalently, for every Borel set \(E\subseteq[0,1]^d\),
\[
\nu_h(E)=h^{-d}\int_{E\cap C_h}dx.
\]
\end{definitionv}

\begin{definitionv}{synth_1}[Scaled coarse polynomial basis]
The coarse polynomial basis vector \(r\), for dimension \(d\in\mathbb N\) and scale \(h\in\mathbb R\), is defined at each \(x\in\mathbb R^d\) by
\[
r(x)=
\left(
\prod_{i=1}^{d}
\begin{cases}
1, & \kappa_i=0,\\
\sqrt{12}\,\dfrac{x_i-\tfrac12}{h}, & \kappa_i=1,\\
\sqrt{5}\left(6\left(\dfrac{x_i-\tfrac12}{h}\right)^2-\tfrac12\right),
& \kappa_i=2
\end{cases}
\right)_{\kappa\in\mathbb N_0^d:\,\sum_{i=1}^{d}\kappa_i\le 2}.
\]
Division by zero is interpreted as zero, so this definition also specifies \(r\) at \(h=0\).
\end{definitionv}

\begin{definitionv}{synth_7}[Localized cellwise polynomial basis]
For \(d\ge1\), \(0<h\le1/2\), and an integer \(J\ge1\), let \(x_0=(1/2,\ldots,1/2)\), \(C_h=x_0+[-h/2,h/2]^d\), and \(\delta=h/J\). Partition each coordinate interval of \(C_h\) into \(J\) intervals of length \(\delta\), taking each interval half-open on the right except for the last interval, which includes the outer right endpoint. Their Cartesian products form the partition \(\{C_{\mathbf j}:\mathbf j\in\{0,\ldots,J-1\}^d\}\). The center of \(C_{\mathbf j}\) is
\[
c_{\mathbf j}=x_0-\frac h2\mathbf1+\delta\left(\mathbf j+\frac12\mathbf1\right).
\]
Write
\[
P_0(u)=1,\qquad P_1(u)=u,\qquad P_2(u)=\frac{3u^2-1}{2}.
\]
For each cell index \(\mathbf j\) and multi-index \(\kappa\in\mathbb N_0^d\) with \(|\kappa|\le2\), define
\[
b_{\mathbf j,\kappa}(x)
=
J^{d/2}\mathbf1\{x\in C_{\mathbf j}\}
\prod_{v=1}^d
\sqrt{2\kappa_v+1}\,
P_{\kappa_v}\!\left(\frac{2(x_v-c_{\mathbf j,v})}{\delta}\right).
\]
The vector \(\mathbf b_J(x)\) collects these functions, ordered lexicographically by cell index and then by multi-index. Its dimension is
\[
k=J^d\binom{d+2}{2}.
\]
These functions form the positive-leading-coefficient cellwise scaled Legendre orthonormal basis, under \(d\nu_h(x)=h^{-d}\mathbf1\{x\in C_h\}\,dx\), of the space \(V_J\) of polynomials of total degree at most two on each partition cell. Set \(\mathbf b_J(x)=0\) for \(x\notin C_h\). Thus the associated kernel is uniquely specified by
\[
K_J(x,x')=\mathbf b_J(x)^\top\mathbf b_J(x').
\]
\end{definitionv}

\begin{definitionv}{def:projection}[Projection \(\Pi_J\) and kernel \(K_J\)]
\[
\delta=\frac hJ,\qquad
K_J(x,x')=\mathbf b_J(x)^\top \mathbf b_J(x'),\qquad
\Pi_Jf(x)=\int K_J(x,x')f(x')\,d\nu_h(x').
\]
Here \(\delta\) is the fine-cell side length, \(h\) is the localization scale, \(J\) is the grid resolution, \(\mathbf b_J\) is the cellwise orthonormal fine polynomial basis vector, and \(\nu_h\) is the uniform probability measure on the localization cube \(C_h\). Weighted basis expressions are zero outside \(C_h\).
\end{definitionv}

\begin{algorithmv}{def:roles}[Record roles \(\mathcal D^T,I_L,I_T\)]
Given the original dataset \(\mathcal D_{n,m}\) of \(n\) labeled records and \(m\) independent outcome-free treatment records, form the two roles as follows.
\begin{enumerate}
\item Concatenate the covariate-treatment pairs, with total treatment-record count \(N\):
\[
N=n+m,\qquad
\mathcal D^T=((X_i^T,A_i^T))_{i=1}^N,
\qquad
(X_i^T,A_i^T)=
\begin{cases}
(X_i,A_i),&1\le i\le n,\\
(\widetilde X_{i-n},\widetilde A_{i-n}),&n<i\le N.
\end{cases}
\]
Here \(X\) denotes the covariate vector and \(A\) the binary treatment indicator; tildes denote auxiliary records.
\item Define the outcome-role count \(\ell\), treatment-role count \(t\), and their index sets:
\[
\ell=\lfloor n/2\rfloor,\qquad
t=N-\ell,\qquad
I_L=\{1,\ldots,\ell\},\qquad
I_T=\{\ell+1,\ldots,N\}.
\]
\item Output the outcome-role labeled records and treatment-role covariate-treatment pairs:
\[
\bigl((X_j,A_j,Y_j)\bigr)_{j\in I_L},
\qquad
\bigl((X_i^T,A_i^T)\bigr)_{i\in I_T}.
\]
Here \(Y_j\) is the observed response. Outcomes from treatment-role records are discarded.
\end{enumerate}
\end{algorithmv}

\begin{algorithmv}{def:moments}[Empirical moments \(\widehat R_{h,J},\widehat Q_{h,J}\)]
Given the original dataset \(\mathcal D_{n,m}\), localization scale \(h\), grid resolution \(J\), localization cube \(C_h\), localization density \(w_h\), coarse polynomial basis vector \(r\), and fine projection kernel \(K_J\), compute the moments as follows.
\begin{enumerate}
\item Form the independent outcome and treatment roles, with counts and index sets
\[
N=n+m,\qquad
\ell=\lfloor n/2\rfloor,\qquad
t=N-\ell,\qquad
I_L=\{1,\ldots,\ell\},\qquad
I_T=\{\ell+1,\ldots,N\}.
\]
The treatment sequence consists of the covariate-treatment pairs
\[
(X_i^T,A_i^T)=
\begin{cases}
(X_i,A_i),&1\le i\le n,\\
(\widetilde X_{i-n},\widetilde A_{i-n}),&n<i\le N.
\end{cases}
\]
Use labeled records \((X_j,A_j,Y_j)\) for \(j\in I_L\) in the outcome role and pairs \((X_i^T,A_i^T)\) for \(i\in I_T\) in the treatment role. Treatment-role outcomes are discarded. Weighted basis expressions are zero outside \(C_h\).
\item Define the outcome-role empirical projection:
\[
\widehat\eta_J^L(x)=
\begin{cases}
\displaystyle\frac1\ell\sum_{j\in I_L}
w_h(X_j)K_J(x,X_j)Y_j,&x\in C_h,\\
0,&x\notin C_h.
\end{cases}
\]
\item Define the empirical response vector:
\[
\begin{aligned}
\widehat R_{h,J}
&=\frac1\ell\sum_{j\in I_L}w_h(X_j)r(X_j)A_jY_j
-\frac1t\sum_{i\in I_T}w_h(X_i^T)r(X_i^T)A_i^T
\widehat\eta_J^L(X_i^T)\\
&=\frac1\ell\sum_{j\in I_L}w_h(X_j)r(X_j)A_jY_j
-\frac1{t\ell}\sum_{i\in I_T}\sum_{j\in I_L}
w_h(X_i^T)w_h(X_j)r(X_i^T)A_i^T
K_J(X_i^T,X_j)Y_j.
\end{aligned}
\]
The response factor in the rectangular correction is \(Y_j\), and the first summand uses \(A_jY_j\). The nested and rectangular expressions define the same vector.
\item Define the unsymmetrized empirical matrix:
\[
\begin{aligned}
\widehat Q^{\mathrm{raw}}_{h,J}
&=\frac1\ell\sum_{j\in I_L}w_h(X_j)r(X_j)r(X_j)^\top A_j\\
&\quad-\frac1{t\ell}\sum_{i\in I_T}\sum_{j\in I_L}
w_h(X_i^T)w_h(X_j)r(X_i^T)r(X_j)^\top
A_i^TA_jK_J(X_i^T,X_j).
\end{aligned}
\]
\item Output the response vector and symmetrized matrix:
\[
\left(\widehat R_{h,J},\widehat Q_{h,J}\right),
\qquad
\widehat Q_{h,J}
=\frac{\widehat Q^{\mathrm{raw}}_{h,J}
+(\widehat Q^{\mathrm{raw}}_{h,J})^\top}{2}.
\]
All inputs are observed in \(\mathcal D_{n,m}\), including when \(m=0\).
\end{enumerate}
\end{algorithmv}

\begin{definitionv}{def:estimator}[Guarded local polynomial decision \(\widehat T_{h,J}\)]
The local polynomial decision \(\widehat T_{h,J}\) is defined by
\[
\widehat T_{h,J}=
\begin{cases}
\mathrm{clip}_{[-1,1]}
\bigl(r_0^\top\widehat Q_{h,J}^{-1}\widehat R_{h,J}\bigr),
&\lambda_{\min}(\widehat Q_{h,J})\ge g(\varepsilon),\\
0,
&\lambda_{\min}(\widehat Q_{h,J})<g(\varepsilon).
\end{cases}
\]
Here \(r_0\) is the coarse polynomial basis vector evaluated at the target profile, \(\widehat Q_{h,J}\) is the symmetrized corrected empirical matrix, and \(\widehat R_{h,J}\) is the corrected empirical response vector. The operator \(\mathrm{clip}_{[-1,1]}\) projects a real value onto \([-1,1]\), and \(\lambda_{\min}\) denotes the smallest eigenvalue. For the public overlap level \(\varepsilon\), the inversion threshold is
\[
g(\varepsilon)=\frac{\varepsilon(1-\varepsilon)}{2}.
\]
\end{definitionv}

\begin{definitionv}{def:tuning}[Public tuning \(h_*,J_*,\widehat T_{\mathrm{up}}\)]
The localization scale \(h_*\), grid resolution \(J_*\), and tuned decision \(\widehat T_{\mathrm{up}}\) are
\[
h_*=\frac12\max\left\{n^{-1/(2\gamma+d)},(nN)^{-1/\Delta}\right\},
\qquad
J_*=\left\lceil h_*^{-\max\{0,\gamma/S-1\}}\right\rceil,
\qquad
\widehat T_{\mathrm{up}}=\widehat T_{h_*,J_*},
\]
where \(N\) is the total treatment-record count, \(S\) is the sum of the propensity and control-mean Hölder orders, and \(\Delta\) is the denominator governing the unequal-channel rate exponent.
\end{definitionv}

\begin{definitionv}{def:sharp-decision}[Sharp-rate decision \(T_*\)]
\[
T_*(\mathcal D_{n,m},U)
=\widehat T_{h_*,J_*}(\mathcal D_{n,m})
=\widehat T_{\mathrm{up}}(\mathcal D_{n,m}).
\]
For every original dataset \(\mathcal D_{n,m}\) and independent uniform randomizer \(U\), this decision uses the prescribed empirical moments, empirical eigenvalue cutoff, clipping, role split, and public tuning parameters \(h_*\) and \(J_*\). Its value is independent of \(U\).
\end{definitionv}

\begin{theoremv}{thm:rectangular-upper}[Rectangular estimation upper bounds]
Fix public parameters satisfying
\begin{itemize}
\item \textbf{(Dimension.)} \(d\) is an integer with \(d\ge1\).
\item \textbf{(Smoothness.)} \(0<\alpha\le1\), \(0<\beta\le1\), and \(1\le\gamma\le3\).
\item \textbf{(Radius.)} \(L>1/2\).
\item \textbf{(Overlap.)} \(0<\varepsilon<1/2\).
\end{itemize}
Let \(\mathcal M\) be the primitive class at these fixed parameters from \cref{obj:def:primitive-class}, and use the absolute-error risk from \cref{obj:def:risk}. There exists a finite constant \(C>0\), depending only on these public parameters, with both of the following guarantees.

For every pair of integers \(n\ge2\) and \(m\ge0\), every \(h\in(0,1/2]\), and every integer \(J\ge1\), the estimator \(\widehat T_{h,J}\) from \cref{obj:def:estimator}, viewed as a decision on the original dataset and randomizer that ignores the randomizer, is Borel measurable and satisfies, for every \(P\in\mathcal M\),
\[
\mathcal R_{n,m}(\widehat T_{h,J},P)
\le
C\min\left\{
1,\,
h^\gamma+
(h/J)^{\alpha+\beta}+
(nh^d)^{-1/2}+
\bigl(n(n+m)h^d(h/J)^d\bigr)^{-1/2}
\right\}.
\]

Define the supplied-propensity point-estimation order and the denominator governing the unequal-channel rate exponent by
\[
r_o(n)=n^{-\gamma/(2\gamma+d)},
\qquad
\Delta=2\gamma+d+\frac{\gamma d}{\alpha+\beta},
\]
and define the upper rate by
\[
r_{\mathrm{up}}(n,m)
=
\max\left\{
r_o(n),\,
\bigl(n(n+m)\bigr)^{-\gamma/\Delta}
\right\}.
\]
For every pair of integers \(n\ge2\) and \(m\ge0\), the publicly tuned decision \(\widehat T_{\mathrm{up}}\) from \cref{obj:def:tuning}, likewise viewed as a decision that ignores the randomizer, is Borel measurable and satisfies, for every \(P\in\mathcal M\),
\[
\mathcal R_{n,m}(\widehat T_{\mathrm{up}},P)
\le C r_{\mathrm{up}}(n,m).
\]
\end{theoremv}

\begin{definitionv}{def:sharp-rate}[Sharp rate \(r_*\)]
The sharp rate at the public class parameters is
\[
r_*(n,m;d,\alpha,\beta,\gamma)
=
\max\left\{
n^{-\gamma/(2\gamma+d)},
[n(n+m)]^{-\gamma/(2\gamma+d+\gamma d/(\alpha+\beta))}
\right\}
=
r_{\mathrm{up}}(n,m),
\]
where \(r_{\mathrm{up}}(n,m)\) is the maximum of the supplied-propensity and unequal-channel rate orders.
\end{definitionv}

\begin{theoremv}{thm:sharp-annotation-frontier}[Sharp annotation frontier]
Fix the covariate dimension \(d\in\mathbb N\), the propensity and control-mean smoothness orders \(\alpha,\beta\), the effect smoothness order \(\gamma\), the Hölder radius \(L\), and the overlap level \(\varepsilon\), satisfying:
\begin{itemize}
\item \textbf{(Dimension.)} \(d\ge1\).
\item \textbf{(Nuisance smoothness.)} \(0<\alpha\le1\) and \(0<\beta\le1\).
\item \textbf{(Effect smoothness.)} \(1\le\gamma\le3\).
\item \textbf{(Radius.)} \(L>1/2\).
\item \textbf{(Overlap.)} \(0<\varepsilon<1/2\).
\end{itemize}
Let \(\mathcal M\) denote the primitive class at these fixed parameters, and let \(r_*\), \(T_*\), and \(R(n,m)\) be the sharp rate, total data-only decision, and minimax risk of \cref{obj:def:sharp-rate,obj:def:sharp-decision,obj:def:minimax}, respectively. There exist real constants \(0<c\le C<\infty\), depending only on the fixed parameters, such that, simultaneously for every \(n,m\in\mathbb N\) with \(n\ge2\), \(T_*\) is Borel measurable and
\[
c\,r_*
\le R(n,m)
\le \sup_{P\in\mathcal M}\mathcal R_{n,m}(T_*,P)
\le C\,r_*.
\]
Moreover, every Borel real-valued decision \(T\) using the original dataset and the independent uniform randomizer satisfies
\[
c\,r_*\le \sup_{P\in\mathcal M}\mathcal R_{n,m}(T,P).
\]

Define the total treatment-record count \(N\), the combined nuisance smoothness \(S\), the exponent denominator \(\Delta\), the smoothness threshold \(S_{\mathrm{crit}}\), the total-record growth exponent \(q_*\), and the supplied-propensity order \(r_o(n)\) by
\[
N=n+m,\qquad S=\alpha+\beta,\qquad
\Delta=2\gamma+d+\frac{\gamma d}{\alpha+\beta},
\]
\[
S_{\mathrm{crit}}=\frac{\gamma d}{2\gamma+d},
\qquad
q_*=\frac{\gamma d}{(\alpha+\beta)(2\gamma+d)},
\qquad
r_o(n)=n^{-\gamma/(2\gamma+d)}.
\]
For these sample sizes, the sharp rate has the branches
\[
r_*=
\begin{cases}
r_o(n), & S\ge S_{\mathrm{crit}},\\
(nN)^{-\gamma/\Delta}, & S<S_{\mathrm{crit}},\quad N\le n^{q_*},\\
r_o(n), & S<S_{\mathrm{crit}},\quad N\ge n^{q_*}.
\end{cases}
\]
Whenever \(N=n^{q_*}\), the two rate formulas coincide:
\[
(nN)^{-\gamma/\Delta}=r_o(n).
\]
\end{theoremv}

\begin{theoremv}{thm:supplied-propensity}[Supplied-propensity point estimation rate]
Fix public parameters satisfying
\begin{itemize}
\item \textbf{(Dimension.)} \(d\in\mathbb N\) and \(d\ge1\).
\item \textbf{(Nuisance smoothness.)} \(0<\alpha\le1\) and \(0<\beta\le1\).
\item \textbf{(Effect smoothness.)} \(1\le\gamma\le3\).
\item \textbf{(Radius.)} \(L>1/2\).
\item \textbf{(Overlap.)} \(0<\varepsilon<1/2\).
\end{itemize}
Let \(\mathcal M\) be the primitive class at these parameters in \cref{obj:def:primitive-class}, and let \(R^e(n,m)\) be the supplied-propensity minimax absolute-error risk in \cref{obj:def:oracle-risk}. Define the supplied-propensity rate and the evaluation profile by
\[
r_o(n)=n^{-\gamma/(2\gamma+d)},
\qquad
x_0=(1/2,\ldots,1/2).
\]
There exist real constants \(0<c\le C<\infty\), depending only on the fixed public parameters, such that, for every \(n,m\in\mathbb N\) with \(n\ge2\),
\[
c\,r_o(n)\le R^e(n,m)\le C\,r_o(n).
\]

The upper bound is attained by the explicit degree-two smoother under the known uniform design, applied to the inverse-propensity pseudo-outcomes
\[
\frac{A_iY_i}{e_P(X_i)}
-
\frac{(1-A_i)Y_i}{1-e_P(X_i)},
\qquad i=1,\ldots,n,
\]
at bandwidth
\[
h=n^{-1/(2\gamma+d)},
\]
with final clipping to \([-1,1]\). Here the smoother uses its localization weight and coarse polynomial basis at this bandwidth for every \(h>0\), including \(h>1/2\). For every \(P\in\mathcal M\), the resulting decision \(T^e\) is measurable as a function of the dataset and randomizer when supplied with the designated admissible propensity \(e_P\), satisfies
\[
\left|T^e(\mathcal D_{n,m},U;e_P)\right|\le1
\]
for every dataset and randomizer value, and obeys
\[
\mathbb E_{\mathsf E_{n,m}(P)}
\left[
\left|T^e(\mathcal D_{n,m},U;e_P)-\tau_P(x_0)\right|
\right]
\le C\,r_o(n).
\]
The sampling law \(\mathsf E_{n,m}(P)\) is the product law of \(n\) independent labeled records, \(m\) independent outcome-free covariate-treatment records, and an independent uniform randomizer \(U\), as specified in \cref{obj:ass:sample-law}.

The same lower bound holds for the supplied-propensity minimax risk over the same-class subfamily with constant propensity and control mean:
\[
c\,r_o(n)\le
\inf_{T^e}
\sup_{\substack{P\in\mathcal M\\ e_P=\mu_{0,P}=1/2}}
\mathbb E_{\mathsf E_{n,m}(P)}
\left[
\left|T^e(\mathcal D_{n,m},U;e_P)-\tau_P(x_0)\right|
\right],
\]
where the constant-witness equalities hold on the unit cube and the infimum ranges over the supplied-propensity decisions of \cref{obj:def:oracle-risk}.
\end{theoremv}

\begin{definitionv}{def:annotation-region}[Annotation region \(\mathcal B_*(\zeta)\)]
The annotation region at tolerance \(\zeta\) is
\[
\mathcal B_*(\zeta)=
\left\{
(n,m)\in\mathbb N_{\ge2}\times\mathbb N_0:
r_*(n,m;d,\alpha,\beta,\gamma)\le\zeta r_o(n)
\right\},
\]
where \(r_*\) is the sharp rate and \(r_o(n)\) is the supplied-propensity point-estimation order.
\end{definitionv}

\begin{lemmav}{lem:local-polynomial-projection-facts}[Local polynomial projection bounds]
Fix public parameters \(d\in\mathbb N\) and \(\alpha,\beta,\gamma,L,\varepsilon\in\mathbb R\) satisfying:
\begin{itemize}
\item \textbf{(Dimension.)} \(d\ge1\).
\item \textbf{(Smoothness.)} \(0<\alpha\le1\), \(0<\beta\le1\), and \(1\le\gamma\le3\).
\item \textbf{(Radius.)} \(L>1/2\).
\item \textbf{(Overlap.)} \(0<\varepsilon<1/2\).
\end{itemize}
There exists a finite constant \(C=C(d,\alpha,\beta,\gamma,L,\varepsilon)>0\) such that, for every primitive law \(P\in\mathcal M(d,\alpha,\beta,\gamma,L,\varepsilon)\) as defined in \cref{obj:def:primitive-class}, every \(h\in\mathbb R\) with \(0<h\le1/2\), and every integer \(J\ge1\), the following conclusions hold for the bases and projection of \cref{obj:def:projection}.

Write \(x_0=(1/2,\ldots,1/2)\) for the evaluation center, \(C_h=[1/2-h/2,1/2+h/2]^d\) for the localization cube, and \(\nu_h\) for its uniform probability measure. Let \(p\) be the Taylor polynomial of the canonical contrast \(\tau_P\) at \(x_0\), with Taylor degree \(p_*=\lceil\gamma\rceil-1\):
\[
p(x)=
\sum_{\substack{\kappa\in\mathbb N^d\\ \sum_{i=1}^d\kappa_i\le p_*}}
\frac{\partial^\kappa\tau_P(x_0)}
     {\prod_{i=1}^d\kappa_i!}
\prod_{i=1}^d(x_i-1/2)^{\kappa_i}.
\]
There exists a coarse coefficient vector \(\vartheta\) whose represented polynomial
\[
p_\vartheta(x)=r(x)^\top\vartheta
\]
satisfies
\[
\sup_{x\in C_h}|\tau_P(x)-p_\vartheta(x)|\le Ch^\gamma,
\qquad
\|\vartheta\|_2\le C,
\qquad
p_\vartheta(x_0)=\tau_P(x_0),
\]
and
\[
p_\vartheta(x)=p(x)\qquad\text{for every }x\in C_h.
\]

The coarse basis vector \(r\) and fine projection kernel \(K_J\) satisfy
\[
\sup_{x\in C_h}\|r(x)\|_2\le C,
\qquad
\sup_{x\in C_h}\int |K_J(x,x')|\,d\nu_h(x')\le C,
\]
and
\[
\iint K_J(x,x')^2\,d\nu_h(x)\,d\nu_h(x')
=qJ^d,
\]
where \(q\) is the number of coarse multi-indices of total degree at most two:
\[
q=\#\left\{\kappa\in\mathbb N^d:\sum_{i=1}^d\kappa_i\le2\right\}.
\]

Finally, with fine-cell side length \(\delta=h/J\), every coarse basis coordinate \(r_u\) satisfies
\[
\|e_Pr_u-\Pi_J(e_Pr_u)\|_{L_2(\nu_h)}
\le C\delta^\alpha,
\]
where \(e_P\) is the designated admissible propensity witness, and the designated admissible control-mean witness \(\mu_{0,P}\) satisfies
\[
\|\mu_{0,P}-\Pi_J\mu_{0,P}\|_{L_2(\nu_h)}
\le C\delta^\beta.
\]
\end{lemmav}

\begin{lemmav}{lem:population-positivity}[Population quadratic form positivity]
Suppose the following conditions hold:
\begin{itemize}
\item \textbf{(Public parameters.)} The dimension \(d\) is an integer with \(d\ge1\), the smoothness parameters satisfy \(0<\alpha,\beta\le1\) and \(1\le\gamma\le3\), the radius satisfies \(L>1/2\), and the overlap parameter satisfies \(0<\varepsilon<1/2\).
\item \textbf{(Primitive law.)} The law \(P\) belongs to the primitive class \(\mathcal M\) at these public parameters, as specified in \cref{obj:def:primitive-class}. Let \(e_P\) be its designated admissible propensity witness.
\item \textbf{(Sample sizes.)} The sample counts \(n,m\) are nonnegative integers with \(n\ge2\).
\item \textbf{(Localization and projection.)} The localization scale satisfies \(0<h\le1/2\), and the grid resolution \(J\) is an integer with \(J\ge1\). Use the projection \(\Pi_J\) from \cref{obj:def:projection}.
\end{itemize}
Let \(\nu_h\) be the uniform localization measure on the localization cube \(C_h\), with
\[
d\nu_h(x)=h^{-d}\mathbf 1_{C_h}(x)\,dx.
\]
For every coefficient vector \(v\in\mathbb R^q\), indexed by the multi-indices of total degree at most two, define its coarse-basis polynomial and the population matrix by
\[
p_v(x)=\sum_u r_u(x)v_u,
\qquad
Q_{P,h,J}
=\mathbb E_{\mathsf E_{n,m}(P)}[\widehat Q_{h,J}],
\]
where \(r_u\) are the scaled orthonormal coarse polynomial basis functions and \(\widehat Q_{h,J}\) is the symmetrized corrected empirical matrix. Then
\[
\begin{aligned}
v^\top Q_{P,h,J}v
&=\int e_P(x)\bigl(1-e_P(x)\bigr)p_v(x)^2\,d\nu_h(x)\\
&\quad+\int
\bigl[e_P(x)p_v(x)-\Pi_J(e_Pp_v)(x)\bigr]^2\,d\nu_h(x)\\
&\ge \varepsilon(1-\varepsilon)\sum_u v_u^2.
\end{aligned}
\]
\end{lemmav}

\begin{lemmav}{lem:rectangular-sampling-bound}[Rectangular sampling bound]
Fix the covariate dimension \(d\in\mathbb N\), the smoothness parameters \(\alpha,\beta,\gamma\in\mathbb R\), the Hölder radius \(L\in\mathbb R\), and the overlap parameter \(\varepsilon\in\mathbb R\), satisfying
\begin{itemize}
\item \textbf{(Dimension.)} \(d\ge1\).
\item \textbf{(Nuisance smoothness.)} \(0<\alpha\le1\) and \(0<\beta\le1\).
\item \textbf{(Effect smoothness.)} \(1\le\gamma\le3\).
\item \textbf{(Radius.)} \(L>1/2\).
\item \textbf{(Overlap parameter.)} \(0<\varepsilon<1/2\).
\end{itemize}
There exists a finite constant \(C=C(d,\alpha,\beta,\gamma,L,\varepsilon)>0\) such that the following bound holds for every primitive law
\(P\in\mathcal M(d,\alpha,\beta,\gamma,L,\varepsilon)\) from \cref{obj:def:primitive-class}, and every \(n,m,J\in\mathbb N\) and \(h\in\mathbb R\) satisfying
\begin{itemize}
\item \textbf{(Sample sizes.)} \(n\ge2\) and \(m\ge0\).
\item \textbf{(Localization.)} \(0<h\le1/2\).
\item \textbf{(Resolution.)} \(J\ge1\).
\end{itemize}
Take expectations under \(\mathsf E_{n,m}(P)\), the product sampling law specified in \cref{obj:ass:sample-law}. Write \(N\) for the total treatment-record count, \(q\) for the number of polynomial multi-indices of total degree at most two, and \(k\) for the fine projection rank:
\[
N=n+m,\qquad q=\binom{d+2}{2},\qquad k=qJ^d.
\]
For the empirical response vector \(\widehat R_{h,J}\) and symmetrized corrected empirical matrix \(\widehat Q_{h,J}\) from \cref{obj:def:moments}, define their population expectations by
\[
\bar R=\mathbb E_{\mathsf E_{n,m}(P)}\widehat R_{h,J},
\qquad
Q_{P,h,J}=\mathbb E_{\mathsf E_{n,m}(P)}\widehat Q_{h,J}.
\]
Then
\[
\mathbb E_{\mathsf E_{n,m}(P)}
 \bigl\|\widehat R_{h,J}-\bar R\bigr\|_2^2
+
\mathbb E_{\mathsf E_{n,m}(P)}
 \bigl\|\widehat Q_{h,J}-Q_{P,h,J}\bigr\|_F^2
\le
C\left\{\frac{1}{nh^d}+\frac{k}{nNh^{2d}}\right\}.
\]
\end{lemmav}

\begin{definitionv}{synth_4}[Treatment-record marginal]
For a primitive Borel probability law
\(P=\mathcal L(X,A,Y(0),Y(1))\) on
\([0,1]^d\times\{0,1\}^3\), define its treatment-record marginal by
\[
P_{XA}=\mathcal L_P(X,A).
\]
Thus \(P_{XA}\) is the probability law on \([0,1]^d\times\{0,1\}\) satisfying
\[
P_{XA}(E)=P\bigl((X,A)\in E\bigr)
\]
for every Borel set \(E\subseteq[0,1]^d\times\{0,1\}\).
\end{definitionv}

\begin{definitionv}{synth_6}[Fixed-size hypothesis mixtures]
For \(n\in\mathbb N\) with \(n\ge2\), \(m\in\mathbb N_0\), and the finite priors \(\varpi_\theta\) produced by \cref{obj:def:marked-handle} on its probability-valued domain, define
\[
M_\theta
:=\int\bigl(P_o^{\otimes n}\otimes P_{XA}^{\otimes m}\bigr)\,d\varpi_\theta(P),
\qquad \theta\in\{-1,+1\},
\]
where \(P_o=\mathcal L_P(X,A,Y)\), \(Y=(1-A)Y(0)+AY(1)\), and \(P_{XA}=\mathcal L_P(X,A)\). Thus \(M_{+1}\) and \(M_{-1}\) are the positive- and negative-hypothesis probability laws of the original-record dataset \(\mathcal D_{n,m}\), with \(n\) labeled records and \(m\) outcome-free treatment records. Under each mixture, one primitive law \(P\) is drawn from the corresponding prior and all \(n+m\) records are drawn conditionally independently from that law through their respective observation channels.
\end{definitionv}

\begin{definitionv}{def:hellinger}[Squared Hellinger distance \(H^2\)]
The squared Hellinger distance between the positive- and negative-hypothesis mixture laws \(M_{+1}\) and \(M_{-1}\) is
\[
H^2(M_{+1},M_{-1})
=
\int
\left(
\sqrt{\frac{dM_{+1}}{d(M_{+1}+M_{-1})}}
-
\sqrt{\frac{dM_{-1}}{d(M_{+1}+M_{-1})}}
\right)^2
\,d(M_{+1}+M_{-1}).
\]
\end{definitionv}

\begin{definitionv}{synth_2}[Constructed arm-mean functions]
The real-valued arm-mean functions \(\mu_{0,\theta,\eta}\) and \(\mu_{1,\theta,\eta}\) are defined for every dimension \(d\in\mathbb N\), real scales \(h,\delta\), real propensity perturbation amplitude \(a\), real response perturbation amplitude \(b\), hypothesis sign \(\theta\in\{-1,+1\}\), and response sign array \(\eta\in\{-1,+1\}^{\mathcal I}\) by
\[
\begin{aligned}
\mu_{0,\theta,\eta}(x)
&=\frac12+b\sum_{\mathbf z\in\mathcal I}
  \eta_{\mathbf z}\phi_{\mathbf z}(x)
  +\theta ab\,B_h(x)^2,\\
\mu_{1,\theta,\eta}(x)
&=\frac12+b\sum_{\mathbf z\in\mathcal I}
  \eta_{\mathbf z}\phi_{\mathbf z}(x)
  -\theta ab\,B_h(x)^2,
\qquad x\in\mathbb R^d.
\end{aligned}
\]
Here \(\mathcal I\), the finite frame index set, \(B_h\), the localization bump, and \(\phi_{\mathbf z}\), the localized frame entries, use the defining formulas in \cref{obj:def:marked-handle}, evaluated at dimension \(d\) and the stated real parameter values. The hypothesis sign and response signs encode Boolean values by assigning \(+1\) to true and \(-1\) to false. Each \(\eta_{\mathbf z}\) is the encoded second component of the Boolean sign pair at \(\mathbf z\); the first component is arbitrary. Control and treatment correspond to false and true arm values, respectively.
\end{definitionv}

\begin{algorithmv}{def:marked-handle}[Finite-prior construction \(\mathfrak C\)]
The construction \(\mathfrak C\) takes inputs \((h,\delta,a,b)\) satisfying \(0<\delta\le h\le1/2\) and \(a,b>0\), at the public parameters \((d,\alpha,\beta,\gamma,L)\) and fixed interior covariate profile \(x_0\).
\begin{enumerate}
\item Define the real functions
\[
\chi(u)=
\begin{cases}
\exp(-1/u),&u>0,\\
0,&u\le0,
\end{cases}
\qquad
\vartheta(u)=\frac{\pi}{2}\frac{\chi(u)}{\chi(u)+\chi(1-u)},
\qquad u\in\mathbb R,
\]
and, for each \(z\in\mathbb Z\),
\[
\omega_z(u)=
\begin{cases}
\sin\vartheta(u-z+1),&z-1\le u\le z,\\
\cos\vartheta(u-z),&z\le u\le z+1,\\
0,&u\notin[z-1,z+1].
\end{cases}
\]
The two expressions at \(u=z\) agree.

\item For \(w\in\mathbb R^d\) and \(x\in[0,1]^d\), define the localization bump and fine-scale frame by
\[
B(w)=\prod_{v=1}^d\frac{\chi(1-4w_v^2)}{\chi(1)},
\qquad
B_h(x)=B\bigl((x-x_0)/h\bigr),
\]
\[
\psi_{\mathbf z}(x)=
\prod_{v=1}^d
\omega_{z_v}\bigl((x_v-x_{0,v})/\delta\bigr),
\]
\[
\mathcal I=
\left\{\mathbf z\in\mathbb Z^d:
|z_v|\le1+\frac{h}{2\delta}
\text{ for every }v\in\{1,\ldots,d\}\right\},
\qquad
\phi_{\mathbf z}(x)=B_h(x)\psi_{\mathbf z}(x),
\quad \mathbf z\in\mathcal I.
\]

\item For sign arrays \(\lambda,\eta\in\{-1,+1\}^{\mathcal I}\), define
\[
F_\lambda(x)=
\sum_{\mathbf z\in\mathcal I}\lambda_{\mathbf z}\phi_{\mathbf z}(x),
\qquad
G_\eta(x)=
\sum_{\mathbf z\in\mathcal I}\eta_{\mathbf z}\phi_{\mathbf z}(x).
\]
For each hypothesis index \(\theta\in\{-1,+1\}\), define the sign-pair correlation and product weights by
\[
\rho_\theta=\frac{\theta}{2},
\qquad
w_\theta(\lambda,\eta)=
\prod_{\mathbf z\in\mathcal I}
\frac{1+\rho_\theta\lambda_{\mathbf z}\eta_{\mathbf z}}{4}.
\]

\item Define the propensity, contrast, and potential-response means by
\[
e_\lambda(x)=\frac12+aF_\lambda(x),
\qquad
\tau_\theta(x)=-4ab\rho_\theta B_h(x)^2,
\]
\[
\mu_{0,\theta,\eta}(x)=
\frac12+bG_\eta(x)-\frac12\tau_\theta(x),
\qquad
\mu_{1,\theta,\eta}(x)=
\frac12+bG_\eta(x)+\frac12\tau_\theta(x).
\]
The probability-valued domain consists of inputs satisfying
\[
e_\lambda(x),\ \mu_{0,\theta,\eta}(x),\
\mu_{1,\theta,\eta}(x)\in[0,1]
\quad
\text{for every }\theta\in\{-1,+1\},\
\lambda,\eta\in\{-1,+1\}^{\mathcal I},\
x\in[0,1]^d.
\]
On this domain, define the primitive Borel law \(P_{\theta,\lambda,\eta}\) by
\[
X\sim\operatorname{Unif}([0,1]^d),
\qquad
\mathcal L\bigl(A,Y(0),Y(1)\mid X=x\bigr)
=
\operatorname{Bernoulli}\bigl(e_\lambda(x)\bigr)
\otimes
\operatorname{Bernoulli}\bigl(\mu_{0,\theta,\eta}(x)\bigr)
\otimes
\operatorname{Bernoulli}\bigl(\mu_{1,\theta,\eta}(x)\bigr).
\]
Here \(X\) is the covariate vector, \(A\) is the binary treatment indicator, and \(Y(0),Y(1)\) are the binary potential responses under control and treatment. The prescribed functions \(e_\lambda,\mu_{0,\theta,\eta},\mu_{1,\theta,\eta}\) are the designated conditional-margin functions of this law.

\item Output the contrasts and finite pushforward measures
\[
\mathfrak C(h,\delta,a,b)
=
\bigl((\tau_\theta,\varpi_\theta)\bigr)_{\theta\in\{-1,+1\}},
\qquad
\varpi_\theta=
\sum_{\lambda,\eta\in\{-1,+1\}^{\mathcal I}}
w_\theta(\lambda,\eta)\,
\operatorname{Dirac}(P_{\theta,\lambda,\eta}),
\]
where \(\operatorname{Dirac}(P)\) is the unit point mass at the primitive law \(P\).
\end{enumerate}
\end{algorithmv}

\begin{lemmav}{lem:hellinger-block-calculus}[Hellinger calculus for blocks]
For probability laws \(\mathbb B,\mathbb C\) on a common measurable space, use the normalization of \cref{obj:def:hellinger}:
\[
\mathsf h^2(\mathbb B,\mathbb C)
=\int\left(\sqrt{\frac{d\mathbb B}{d\xi}}
           -\sqrt{\frac{d\mathbb C}{d\xi}}\right)^2\,d\xi,
\qquad \xi=\mathbb B+\mathbb C.
\]
The following statements hold on arbitrary measurable spaces.

For any finite family of pairs of probability laws \(\mathbb B_{u_*},\mathbb C_{u_*}\), each pair on its respective common measurable space,
\[
\mathsf h^2\left(\bigotimes_{u_*}\mathbb B_{u_*},
                \bigotimes_{u_*}\mathbb C_{u_*}\right)
\le \sum_{u_*}\mathsf h^2(\mathbb B_{u_*},\mathbb C_{u_*}).
\]

Let \(\nu\) be a probability law on the conditioning space, and let \(\mathbb B_z,\mathbb C_z\) be probability kernels into a common measurable space. Suppose there exists a common s-finite reference kernel such that both kernels are represented, at every \(z\), by nonnegative extended-real densities with respect to that reference kernel, with both densities jointly measurable in the conditioning and outcome coordinates. Then
\[
\mathsf h^2\bigl(\nu(dz)\mathbb B_z(dw),
                \nu(dz)\mathbb C_z(dw)\bigr)
=\int \mathsf h^2(\mathbb B_z,\mathbb C_z)\,d\nu(z).
\]

For any probability laws \(\mathbb B,\mathbb C\) on a common measurable space and any probability law \(\nu\) on another measurable space,
\[
\mathsf h^2(\mathbb B\otimes\nu,\mathbb C\otimes\nu)
=\mathsf h^2(\mathbb B,\mathbb C).
\]
\end{lemmav}

\begin{lemmav}{lem:published-fuzzy-testing}[Absolute-loss fuzzy testing bound]
Let \((\mathbb B_z)\) and \((\mathbb C_z)\) be probability kernels from a common measurable index space to a common measurable observation space, and let \(\omega\) be a probability prior on the index space. Let \(v_0\ge1\) be an integer, let \(\mathcal M\) be a model of measures containing \(\mathbb B_z\) and \(\mathbb C_z\) for every \(z\), and let \(\Psi\) be a real-valued functional on measures. Use the squared Hellinger distance of \cref{obj:def:hellinger}, with normalization
\[
\mathsf h^2(\mathbb B,\mathbb C)
=
\int
\left(
\sqrt{\frac{d\mathbb B}{d\xi}}
-
\sqrt{\frac{d\mathbb C}{d\xi}}
\right)^2\,d\xi,
\qquad
\xi=\mathbb B+\mathbb C.
\]
Suppose the following conditions hold:
\begin{itemize}
\item \textbf{(Constants.)} The separation bound \(s_0\) and distance bound \(u_0\) satisfy \(s_0>0\) and \(0\le u_0<2\).
\item \textbf{(Separation.)} For every pair of indices \(z,z'\),
\[
\Psi(\mathbb B_z)-\Psi(\mathbb C_{z'})\ge s_0.
\]
\item \textbf{(Mixture distance.)} The product-law mixtures satisfy
\[
\mathsf h^2\left(
\int \mathbb B_z^{\otimes v_0}\,d\omega(z),
\int \mathbb C_z^{\otimes v_0}\,d\omega(z)
\right)\le u_0.
\]
\end{itemize}
Then every measurable real-valued decision \(\widehat\Psi\) on the \(v_0\)-observation sample space satisfies
\[
\sup_{P\in\mathcal M}
\int
\left|\widehat\Psi(w)-\Psi(P)\right|
\,dP^{\otimes v_0}(w)
\ge
\frac{s_0}{4}
\left\{1-\sqrt{u_0(1-u_0/4)}\right\},
\]
where the nonnegative loss integral is interpreted in the extended sense. This is the absolute-loss specialization of \citep[Section 3, Lemma 1]{KennedyBalakrishnanRobinsWasserman2024}.
\end{lemmav}

\begin{theoremv}{thm:marked-component-certificate}[Marked component mixture certificate]
Fix the primitive class \(\mathcal M\) of \cref{obj:def:primitive-class}, with public parameters satisfying
\[
d\in\mathbb N,\qquad d\ge1,\qquad
\alpha,\beta\in(0,1],\qquad
\gamma\in[1,3],\qquad
L>\tfrac12,\qquad
\varepsilon\in(0,\tfrac12).
\]
There exist constants \(c\in(0,1]\), \(c_0>0\), and \(C>0\), depending only on these public parameters, such that the following holds for every choice satisfying:
\begin{itemize}
\item \textbf{(Record counts.)} \(n,m\in\mathbb N\), \(n\ge2\), with total record count \(N=n+m\).
\item \textbf{(Scales.)} \(0<\delta\le h\le\tfrac12\).
\item \textbf{(Propensity amplitude.)} \(0<a\le c\delta^\alpha\).
\item \textbf{(Response amplitude.)} \(0<b\le c\delta^\beta\).
\item \textbf{(Effect calibration.)} \(ab\le ch^\gamma\).
\item \textbf{(Occupancy.)} \(N\delta^d\le c_0\).
\end{itemize}
The prescribed construction \(\mathfrak C\) of \cref{obj:def:marked-handle}, evaluated at \(d,h,\delta,a,b\), has the following properties for both hypothesis indices \(\theta\in\{-1,+1\}\).

Its weights on the finite sign arrays are nonnegative and normalized:
\[
w_\theta(\lambda,\eta)\ge0
\quad\text{for every }\lambda,\eta\in\{-1,+1\}^{\mathcal I},
\qquad
\sum_{\lambda,\eta\in\{-1,+1\}^{\mathcal I}}
w_\theta(\lambda,\eta)=1.
\]
Every constructed law \(P_{\theta,\lambda,\eta}\) belongs to \(\mathcal M\), and its canonical contrast agrees with the deterministic constructed contrast everywhere on the unit cube:
\[
\tau_{P_{\theta,\lambda,\eta}}(x)=\tau_\theta(x)
\qquad
\bigl(x\in[0,1]^d\bigr).
\]
Thus the finite prior \(\varpi_\theta\) is supported on \(\mathcal M\). The designated propensity functions, restricted to the unit cube, have the same distribution under \(\varpi_{+1}\) and \(\varpi_{-1}\).

At the interior center \(x_0=(1/2,\ldots,1/2)\), the contrasts satisfy
\[
\tau_\theta(x_0)=-2\theta ab,
\qquad
\tau_{-1}(x_0)-\tau_{+1}(x_0)=4ab,
\qquad
cab\le4ab.
\]
For the fixed-size original-record mixture laws \(M_{+1}\) and \(M_{-1}\), the unhalved squared Hellinger distance of \cref{obj:def:hellinger} satisfies
\[
H^2(M_{+1},M_{-1})
\le CnNh^d\delta^da^2b^2.
\]

For every integer \(k_{\mathrm{aux}}\ge0\), the auxiliary-record mixtures agree:
\[
\int P_{XA}^{\otimes k_{\mathrm{aux}}}\,d\varpi_{+1}(P)
=
\int P_{XA}^{\otimes k_{\mathrm{aux}}}\,d\varpi_{-1}(P).
\]
Their explicit conditional auxiliary kernels also agree at every collection of \(k_{\mathrm{aux}}\) covariates in \([0,1]^d\). For either hypothesis and every \(x\in[0,1]^d\), the singleton labeled mixture conditional on \(X=x\) is uniform on the four values of \((A,Y)\in\{0,1\}^2\).

Conditional on every collection of all \(N\) covariates in the unit cube, form the graph whose vertices are all \(N\) records. Join two distinct vertices when both covariates lie in the localization cube \(C_h\) and their maximum-norm distance is at most \(2\delta\); records exterior to \(C_h\) are isolated vertices. The conditional mixed density equals the product of the component densities over all connected components of this graph, including the exterior singleton components. Every auxiliary-only component and every component consisting of a single labeled record has the same conditional density, and hence the same conditional law, under the two hypotheses.
\end{theoremv}

\begin{lemmav}{lem:nuisance-frontier-converse}[Nuisance frontier lower bound]
Fix \(d\in\mathbb N\) and \(\alpha,\beta,\gamma,L,\varepsilon\in\mathbb R\), where \(\varepsilon\) is the overlap level. Define the nuisance smoothness sum \(S\), the smoothness threshold \(S_{\mathrm{crit}}\), the record-growth exponent \(q_*\), and the rate denominator \(\Delta\) by
\[
S=\alpha+\beta,\qquad
S_{\mathrm{crit}}=\frac{\gamma d}{2\gamma+d},\qquad
q_*=\frac{\gamma d}{(\alpha+\beta)(2\gamma+d)},\qquad
\Delta=2\gamma+d+\frac{\gamma d}{\alpha+\beta}.
\]
Suppose that:
\begin{itemize}
\item \textbf{(Dimension.)} \(d\ge1\).
\item \textbf{(Nuisance smoothness.)} \(0<\alpha\le1\) and \(0<\beta\le1\).
\item \textbf{(Effect smoothness.)} \(1\le\gamma\le3\).
\item \textbf{(Radius.)} \(L>1/2\).
\item \textbf{(Overlap level.)} \(0<\varepsilon<1/2\).
\item \textbf{(Frontier regime.)} \(S<S_{\mathrm{crit}}\).
\end{itemize}
There exists a constant \(c>0\), depending on the fixed public parameters \(d,\alpha,\beta,\gamma,L,\varepsilon\), such that for every \(n,m\in\mathbb N\) with \(n\ge2\), the total treatment-record count
\[
N=n+m
\]
satisfying \(N\le n^{q_*}\) obeys the following bound for the minimax absolute-error risk \(R(n,m)\) of \cref{obj:def:minimax}, evaluated at these fixed parameters:
\[
R(n,m)\ge c(nN)^{-\gamma/\Delta}.
\]
\end{lemmav}

\begin{propositionv}{prop:annotation-threshold}[Annotation thresholds and acquisition requirements]
Fix public parameters satisfying
\begin{itemize}
\item \textbf{(Dimension.)} \(d\in\mathbb N\) and \(d\ge1\).
\item \textbf{(Nuisance smoothness.)} \(0<\alpha\le1\) and \(0<\beta\le1\).
\item \textbf{(Effect smoothness.)} \(1\le\gamma\le3\).
\item \textbf{(Radius.)} \(L>1/2\).
\item \textbf{(Overlap.)} \(0<\varepsilon<1/2\).
\end{itemize}
Let \(R(n,m)\) and \(R^e(n,m)\) be the minimax and supplied-propensity risks of \cref{obj:def:minimax,obj:def:oracle-risk}, and let \(r_*(n,m)\) be the sharp rate of \cref{obj:def:sharp-rate}, with the fixed parameter arguments suppressed. Define the total treatment-record count \(N\), the summed nuisance smoothness \(S\), the tuning denominator \(\Delta\), the smoothness cutoff \(S_{\mathrm{crit}}\), the total-record exponent \(q_*\), and the oracle rate \(r_o\) by
\[
N=n+m,\qquad S=\alpha+\beta,\qquad
\Delta=2\gamma+d+\frac{\gamma d}{S},
\qquad
S_{\mathrm{crit}}=\frac{\gamma d}{2\gamma+d},
\qquad
q_*=\frac{\gamma d}{S(2\gamma+d)},
\qquad
r_o(n)=n^{-\gamma/(2\gamma+d)}.
\]
All constants below may depend on the fixed public parameters.

For every sequence \(m_n\in\mathbb N\), the following equivalences hold:
\[
\begin{aligned}
\frac{R(n,m_n)}{R^e(n,m_n)}=O(1)
&\quad\Longleftrightarrow\quad
r_*(n,m_n)=O(r_o(n))\\
&\quad\Longleftrightarrow\quad
\exists\,\zeta\in[1,\infty)\\ 
\forall\,n\ \text{ sufficiently large},\\
(n,m_n)\in\mathcal B_*(\zeta)\\
&\quad\Longleftrightarrow\quad
\liminf_{n\to\infty}\frac{n+m_n}{n^{q_*}}>0,
\end{aligned}
\]
where \(\mathcal B_*(\zeta)\) is the region of \cref{obj:def:annotation-region}; the lower limit is understood in the extended nonnegative reals. Moreover,
\[
\begin{aligned}
\frac{R(n,m_n)}{R(n,0)}\longrightarrow0
&\quad\Longleftrightarrow\quad
\frac{r_*(n,m_n)}{r_*(n,0)}\longrightarrow0\\
&\quad\Longleftrightarrow\quad
S<S_{\mathrm{crit}}
\ \text{ and }\\ 
\frac{m_n}{n}\longrightarrow\infty.
\end{aligned}
\]
All sequence limits and order relations are as \(n\to\infty\). If \(S<S_{\mathrm{crit}}\), then \(q_*>1\). For every sequence with \(m_n/n=O(1)\),
\[
R(n,m_n)=O(R(n,0)),
\qquad
R(n,0)=O(R(n,m_n)).
\]
If \(S\ge S_{\mathrm{crit}}\), then every auxiliary-size sequence satisfies
\[
\frac{R(n,m_n)}{R^e(n,m_n)}=O(1).
\]

There exist real constants \(c,C\) with \(0<c\le C\) such that, for every \(n,m\in\mathbb N\) with \(n\ge2\),
\[
c\,r_*(n,m)\le R(n,m)\le C\,r_*(n,m).
\]
For these sample sizes, the exact rate branches are
\[
r_*(n,m)=
\begin{cases}
r_o(n),&S\ge S_{\mathrm{crit}},\\
(nN)^{-\gamma/\Delta},
&S<S_{\mathrm{crit}},\ N\le n^{q_*},\\
r_o(n),
&S<S_{\mathrm{crit}},\ N\ge n^{q_*}.
\end{cases}
\]
At every boundary pair with \(N=n^{q_*}\),
\[
(nN)^{-\gamma/\Delta}=r_o(n).
\]
In particular, for every \(n\ge2\),
\[
r_*(n,0)=
\begin{cases}
n^{-2\gamma/\Delta},&S<S_{\mathrm{crit}},\\
r_o(n),&S\ge S_{\mathrm{crit}}.
\end{cases}
\]
Writing \(s_{\mathrm{pub}}=S/2\) for the average nuisance smoothness, the numerical supervised benchmark coincides with \(r_*(n,0)\): for every \(n\in\mathbb N\),
\[
\begin{cases}
n^{-1/(1+d/(2\gamma)+d/(4s_{\mathrm{pub}}))},
&s_{\mathrm{pub}}<\dfrac{d/4}{1+d/(2\gamma)},\\
n^{-1/(2+d/\gamma)},
&s_{\mathrm{pub}}\ge\dfrac{d/4}{1+d/(2\gamma)}
\end{cases}
=r_*(n,0).
\]
In this identity at \(n=0\), negative powers of zero are assigned the value zero.

Finally, there exist real constants \(c,C>0\), separate from the preceding pair, such that the following acquisition guarantees hold for every target absolute-error scale \(0<\eta<1\) and every \(n,m\in\mathbb N\) with \(n\ge2\). Let \(\mathcal M\) be the primitive class at the fixed parameters from \cref{obj:def:primitive-class}, and use the risk of \cref{obj:def:risk}.

If there exists a Borel decision \(T\) satisfying
\[
\mathcal R_{n,m}(T,P)\le\eta
\qquad\text{for every }P\in\mathcal M,
\]
then
\[
c\,\eta^{-(2+d/\gamma)}\le n,
\qquad
c\,\eta^{-(2+d/\gamma+d/S)}\le n(n+m).
\]
Conversely, if
\[
C\,\eta^{-(2+d/\gamma)}\le n,
\qquad
C\,\eta^{-(2+d/\gamma+d/S)}\le n(n+m),
\]
then the sharp decision \(T_*\) of \cref{obj:def:sharp-decision} is Borel and satisfies
\[
\mathcal R_{n,m}(T_*,P)\le\eta
\qquad\text{for every }P\in\mathcal M.
\]
The exact rate-level acquisition condition is
\[
r_*(n,m)\le\eta
\quad\Longleftrightarrow\quad
\eta^{-(2+d/\gamma)}\le n
\ \text{ and }\
\eta^{-(2+d/\gamma+d/S)}\le n(n+m).
\]
The total treatment-record count also satisfies \(n+m\ge n\).
\end{propositionv}

\begin{lemmav}{lem:published-cate-benchmark}[Published CATE benchmark identities]
Let the dimension \(d\in\mathbb N\) and smoothness orders \(\alpha,\beta,\gamma\in\mathbb R\) satisfy:
\begin{itemize}
\item \textbf{(Dimension.)} \(1\le d\).
\item \textbf{(First nuisance smoothness.)} \(0<\alpha\).
\item \textbf{(Second nuisance smoothness.)} \(0<\beta\).
\item \textbf{(Effect smoothness.)} \(1\le\gamma\).
\end{itemize}
Put \(s_{\mathrm{pub}}=(\alpha+\beta)/2\), the average nuisance smoothness. The numerical exponents and smoothness cutoff displayed in \citet[Section 3, Theorem 1]{KennedyBalakrishnanRobinsWasserman2024} satisfy
\[
\frac{1}{1+d/(2\gamma)+d/(4s_{\mathrm{pub}})}
=\frac{2\gamma}{2\gamma+d+\gamma d/(\alpha+\beta)},
\qquad
\frac{1}{2+d/\gamma}=\frac{\gamma}{2\gamma+d},
\]
and
\[
s_{\mathrm{pub}}<\frac{d/4}{1+d/(2\gamma)}
\quad\Longleftrightarrow\quad
\alpha+\beta<\frac{\gamma d}{2\gamma+d}.
\]
For every \(n\in\mathbb N\) with \(n\ge2\), the numerical piecewise benchmark satisfies
\[
\begin{cases}
n^{-1/(1+d/(2\gamma)+d/(4s_{\mathrm{pub}}))},
& s_{\mathrm{pub}}<(d/4)/(1+d/(2\gamma)),\\
n^{-1/(2+d/\gamma)},
& s_{\mathrm{pub}}\ge(d/4)/(1+d/(2\gamma))
\end{cases}
=
\max\left\{
n^{-\gamma/(2\gamma+d)},
n^{-2\gamma/(2\gamma+d+\gamma d/(\alpha+\beta))}
\right\}.
\]
\end{lemmav}

\begin{definitionv}{def:frontier-handle}[Rate and decision handle]
The frontier handle is the ordered pair \((r_*,T_*)\), parameterized by \(d,n,m\in\mathbb N\) and \(\alpha,\beta,\gamma,\varepsilon\in\mathbb R\), where \(\varepsilon\) is the overlap parameter. Its components are the sharp rate and decision of \cref{obj:def:sharp-rate,obj:def:sharp-decision}:
\[
r_*=
\max\left\{
n^{-\gamma/(2\gamma+d)},
\bigl(n(n+m)\bigr)^{-\gamma/(2\gamma+d+\gamma d/(\alpha+\beta))}
\right\},
\qquad
T_*(\mathcal D_{n,m},U)
=
\widehat T_{\mathrm{up}}(\mathcal D_{n,m}).
\]
Here \(\mathcal D_{n,m}\) is the original dataset and \(U\in\mathbb R\) is the randomizer coordinate. The decision component uses the publicly tuned estimator of \cref{obj:def:tuning} with parameters \(d,\alpha,\beta,\gamma,\varepsilon,n,m\).
\end{definitionv}
